\documentclass[11pt]{article}
\usepackage{explorer}
\newcommand{\ExplorerLanguage}{FR}
\newcommand{\ExplorerNumber}{12}
\newcommand{\ExplorerKind}{Exercices progressifs}
\begin{document}
\ExplorerTitle{Terminale · analyse et algorithmique}{Encadrer une racine par dichotomie}
\ExplorerSection{Données}
Pour l’exercice 1 : on cherche un zéro de \(f(x)=x^{2}-2\) sur [1 ; 2]. Après deux partages, l’intervalle conservé est [1,25 ; 1,5].\par
\vspace{5pt}
\ExplorerSection{Exercice 1}
Poursuivre l’exemple une troisième fois.\par
\vspace{5pt}
\ExplorerSection{Exercice 2}
Partant d’un intervalle de largeur 1, combien de partages garantissent une largeur au plus \(\frac{1}{1000}\) ?\par
\vspace{5pt}
\ExplorerSection{Exercice 3}
Pourquoi le changement de signe de \(f(x)=\frac{1}{x}\) entre −1 et 1 ne prouve-t-il pas l’existence d’un zéro ?\par
\vspace{5pt}
\ExplorerSection{Exercice 4}
Une racine est dans [1,4 ; 1,5]. Quelle erreur maximale garantit le milieu 1,45 ?\par
\vspace{5pt}
\end{document}
